\section{Project Overview}
\label{section:intro}

\subsection{Profile}

\paragraph{Track: 3. Application}

\paragraph{Abstract and project goal:}
Hyperparameter optimization (HPO) is a fundamental yet costly challenge in machine learning: finding the right hyperparameter configuration often requires evaluating dozens of candidate configurations, each demanding a full or partial training run. Reinforcement learning (RL)-based HPO methods, such as Hyp-RL, have shown promise by learning a sequential decision policy over the hyperparameter space. However, existing RL-based approaches treat each configuration evaluation as a black-box query, discarding the rich temporal structure embedded in the learning curve produced during training.

In this project, we propose a learning curve-informed RL framework for HPO and explore two directions. (i) Derivative-informed state representation: We augment the agent's input at each step with explicit first- and second-order derivatives of the observed learning curve, providing structured inductive bias about whether a configuration is still improving, decelerating, or plateauing, at zero additional evaluation cost. (ii) Cross-dataset policy transfer: We investigate whether pre-training the LSTM policy across multiple datasets can yield an agent that generalizes to unseen datasets without restarting from scratch.

\paragraph{TL;DR (``Too Long; Didn't Read"):}
We augment an RL-based HPO agent with learning curve derivatives as structured state features, and investigate cross-dataset transfer of the learned policy for more efficient and generalizable hyperparameter search.

\subsection{Motivation}

\begin{itemize}
    \item \textbf{Why is the problem interesting?}\\
    When human experts tune hyperparameters, they rely on the geometry of the learning curve---observing how fast the loss is dropping or whether it is plateauing---rather than treating each evaluation as an isolated scalar. We aim to provide exactly this information to an RL agent and examine whether it improves tuning efficiency. Beyond single-task HPO, a practically useful agent should also transfer its experience to new datasets, exploiting the structural regularity that learning curves share across tasks.

    \item \textbf{Critical challenges.}\\
    Challenge 1: Reliable derivative estimation from noisy learning curves.
    Learning curves in practice are noisy, and finite-difference estimates can amplify this noise, potentially misleading the agent. Key design questions include window size, smoothing strategy, and sensitivity of the learned policy to these choices---all of which require empirical ablation.\\
    Challenge 2: Learning a transferable HPO policy across datasets.
    A practically useful HPO agent should generalize across tasks---given experience on a collection of datasets, it should be able to search effectively on a new, unseen dataset without retraining from scratch. This is non-trivial because different datasets induce different reward scales, convergence speeds, and hyperparameter sensitivities, making it difficult to learn representations that remain meaningful across tasks.

    \item \textbf{Justify why the problem remains open or unsolved.}\\
    We hypothesize that these derivative features provide a universal view of training dynamics. Since their physical meanings don't change between tasks, they make it easier for the RL agent to transfer its knowledge to new problems. Existing RL-based HPO works observe only scalar rewards, ignoring within-curve structure \citep{jomaa2019hyp}. Learning curve-aware methods operate outside the RL framework \citep{wistuba2022dyhpo}. Meta-RL approaches for HPO address cross-task generalization but rely on complex encoders without exploiting derivative structure \citep{albrechts2024model, wu2023hyperparameter}. Our work is the first to combine derivative-informed state representation with cross-dataset transfer in a unified RL-for-HPO framework.

    \item \textbf{State-of-the-art methods.}\\
    We consider Hyp-RL \citep{jomaa2019hyp} as the state-of-the-art RL-based HPO baseline that our method directly builds upon.
\end{itemize}